% =====================================================================
%  BÁO CÁO THỰC HÀNH — LAB 01: SQL TRÊN GOOGLE CLOUD SHELL VỚI SQLITE
%  Introduction to Data Science — PTIT
%
%  Biên dịch (pdfLaTeX, khuyến nghị, hỗ trợ tiếng Việt tốt nhất):
%      pdflatex lab01_report.tex
%      pdflatex lab01_report.tex      % chạy 2 lần để có mục lục + tham chiếu
%  hoặc:  latexmk -pdf lab01_report.tex
%  hoặc:  bash compile.sh
%
%  Trên Overleaf: upload thư mục này, chọn compiler pdfLaTeX, bấm Recompile.
%  Cần các gói TeX Live: vntex (đi kèm babel-vietnamese), booktabs, listings, tikz,
%  longtable, tabularx, multirow, enumitem, caption, fancyhdr, hyperref.
%  Nếu máy không có vntex, xem mục "Phương án XeLaTeX" trong report/README.md.
% =====================================================================
\documentclass[12pt,a4paper]{article}

% ---------- Tiếng Việt (pdfLaTeX + T5 + vntex/babel) ----------
\usepackage[utf8]{inputenc}
\usepackage[T5]{fontenc}
\usepackage[vietnamese]{babel}

% ---------- Trang, font, bảng biểu ----------
\usepackage[a4paper,top=2.5cm,bottom=2.5cm,left=2.8cm,right=2.5cm,headheight=15pt]{geometry}
\usepackage{amsmath,amssymb}
\usepackage{graphicx}
\usepackage{xcolor}
\usepackage{booktabs}
\usepackage{tabularx}
\usepackage{longtable}
\usepackage{array}
\usepackage{multirow}
\usepackage{enumitem}
\usepackage{float}
\usepackage{caption}
\usepackage{listings}
\usepackage{tikz}
\usetikzlibrary{arrows.meta,positioning}
\usepackage{fancyhdr}
\usepackage[hidelinks,unicode]{hyperref}

% ---------- Công tắc phụ lục ----------
% \withappendixtrue  : bản đầy đủ (báo cáo + phụ lục mã nguồn, ảnh minh chứng)
% \withappendixfalse : bản gọn ~3 trang để nộp nếu giảng viên yêu cầu đúng 2–3 trang
\newif\ifwithappendix
\withappendixtrue

% ---------- Định dạng ----------
\setlength{\parskip}{4pt}
\setlength{\parindent}{0pt}
\renewcommand{\arraystretch}{1.15}
\captionsetup{font=small,labelfont=bf,skip=6pt}

% Tên tiếng Việt cho các thành phần tự động
\renewcommand{\contentsname}{Mục lục}
\renewcommand{\refname}{Tài liệu tham khảo}
\renewcommand{\figurename}{Hình}
\renewcommand{\tablename}{Bảng}
\renewcommand{\appendixname}{Phụ lục}
\renewcommand{\lstlistingname}{Đoạn mã}

% Lệnh tắt cho mã trong dòng chữ — nhớ escape _ % & # khi dùng
\newcommand{\code}[1]{\texttt{\small #1}}

\definecolor{codebg}{gray}{0.96}
\definecolor{codeframe}{gray}{0.72}
\definecolor{codekw}{RGB}{0,80,150}
\definecolor{codecm}{RGB}{0,110,60}

\lstset{
  basicstyle=\ttfamily\footnotesize,
  backgroundcolor=\color{codebg},
  frame=single,
  rulecolor=\color{codeframe},
  keywordstyle=\color{codekw}\bfseries,
  commentstyle=\color{codecm}\itshape,
  stringstyle=\color{red!60!black},
  showstringspaces=false,
  breaklines=true,
  columns=fullflexible,
  keepspaces=true,
  tabsize=2,
  aboveskip=6pt,
  belowskip=6pt,
  captionpos=b
}

\pagestyle{fancy}
\fancyhf{}
\fancyhead[L]{\small Introduction to Data Science}
\fancyhead[R]{\small LAB 01 --- SQL với SQLite}
\fancyfoot[C]{\thepage}
\renewcommand{\headrulewidth}{0.3pt}

\hypersetup{
  pdftitle={Báo cáo LAB 01 - SQL trên Google Cloud Shell với SQLite},
  pdfsubject={Introduction to Data Science - PTIT},
  pdfkeywords={SQL, SQLite, index, EXPLAIN QUERY PLAN, Cloud Shell}
}

\begin{document}

% =====================================================================
%  TRANG BÌA
% =====================================================================
\begin{titlepage}
\centering
\vspace*{1.0cm}

{\large TRƯỜNG ĐẠI HỌC BƯU CHÍNH VIỄN THÔNG (PTIT)\par}
\vspace{0.2cm}
{\large KHOA CÔNG NGHỆ THÔNG TIN\par}

\vspace{1.6cm}
{\color{blue!45!black}\rule{\linewidth}{0.6pt}}\\[0.7cm]
{\Huge\bfseries BÁO CÁO THỰC HÀNH\par}
\vspace{0.6cm}
{\LARGE\bfseries LAB 01\par}
\vspace{0.4cm}
{\Large SQL trên Google Cloud Shell với SQLite\par}
\vspace{0.5cm}
{\large Thiết kế cơ sở dữ liệu quan hệ, CRUD, JOIN, tổng hợp,\par}
{\large subquery, index và phân tích kế hoạch thực thi\par}
\vspace{0.7cm}
{\color{blue!45!black}\rule{\linewidth}{0.6pt}}

\vspace{1.8cm}
\begin{minipage}{0.78\linewidth}
\large
\begin{tabular}{@{}ll@{}}
\textbf{Môn học}            & Introduction to Data Science \\
\textbf{Giảng viên}         & \dotfill \\
\textbf{Sinh viên thực hiện}& \dotfill \\
\textbf{Mã số sinh viên}    & \dotfill \\
\textbf{Lớp}                & \dotfill \\
\textbf{Nhóm}               & \dotfill \\
\end{tabular}
\end{minipage}

\vfill
{\large Hà Nội, \today\par}
\end{titlepage}

% =====================================================================
%  TÓM TẮT
% =====================================================================
\section*{Tóm tắt}
\addcontentsline{toc}{section}{Tóm tắt}

Báo cáo trình bày kết quả thực hành \textbf{LAB 01} của môn Introduction to Data Science: xây
dựng một cơ sở dữ liệu quan hệ bán hàng trực tuyến gồm bốn bảng
(\code{customers}, \code{products}, \code{orders}, \code{order\_items}) trên SQLite, thực hiện
thành thạo các nhóm thao tác của SQL, rồi dùng chỉ mục (index) cùng
\code{EXPLAIN QUERY PLAN} để chứng minh và giải thích ảnh hưởng của index đến kế hoạch
thực thi.

Toàn bộ 21 yêu cầu đã được cài đặt và \textbf{chạy kiểm chứng thực tế}: 15 câu bắt buộc
(S01--S15) chạy trên dữ liệu mẫu của đề, 6 câu bổ sung/về nhà (S16--S21) chạy trên bộ dữ liệu
mở rộng 100 khách hàng -- 30 sản phẩm -- 500 đơn hàng -- 1.500 dòng đơn hàng. Kết quả chính:
mọi câu lệnh chạy không lỗi, dữ liệu luôn thoả ràng buộc khoá ngoại, kế hoạch thực thi chuyển
từ \code{SCAN orders} sang \code{SEARCH orders USING COVERING INDEX} sau khi tạo index trên
\code{orders.customer\_id}.

\vspace{0.4cm}
\noindent\textbf{Từ khoá:} SQLite, SQL, khoá ngoại, JOIN, GROUP BY, HAVING, subquery,
window function, index, EXPLAIN QUERY PLAN, Google Cloud Shell.

\clearpage
\tableofcontents
\clearpage

% =====================================================================
\section{Giới thiệu và mục tiêu}
% =====================================================================

\subsection{Bối cảnh}

Bài thực hành mô phỏng một cửa hàng trực tuyến nhỏ. Dữ liệu được chuẩn hoá thành bốn bảng quan
hệ, cho phép trả lời các câu hỏi kinh doanh quen thuộc: đơn nào giá trị cao nhất, khách nào chi
tiêu nhiều nhất, sản phẩm nào bán chạy, tháng nào doanh thu tăng. Việc thực hành trên SQLite
giúp tiếp cận SQL mà không cần cài đặt máy chủ và không phát sinh chi phí, đúng tinh thần của
đề bài: chỉ dùng Gmail, trình duyệt và Google Cloud Shell.

\subsection{Mục tiêu cụ thể}

\begin{itemize}[leftmargin=1.2cm,itemsep=2pt]
  \item Tạo lập cơ sở dữ liệu quan hệ có khoá chính, khoá ngoại và ràng buộc \code{CHECK};
  \item Thực hiện CRUD (thêm, đọc, sửa, xoá) đúng ngữ nghĩa và an toàn dữ liệu;
  \item Viết truy vấn \code{JOIN}, \code{GROUP BY}, \code{HAVING}, subquery và window function;
  \item Tạo index, đọc và giải thích kế hoạch thực thi do \code{EXPLAIN QUERY PLAN} trả về;
  \item Sinh dữ liệu giả lập quy mô lớn hơn để so sánh hành vi truy vấn khi dữ liệu tăng.
\end{itemize}

\subsection{Phân chia phần bắt buộc và phần bổ sung}

Đề bài có hai khối yêu cầu với \emph{điều kiện dữ liệu khác nhau}, vì vậy bài làm tách thành hai
file đáp án riêng biệt (không trộn lẫn), giúp người chấm thấy ngay đâu là yêu cầu tối thiểu và
đâu là phần mở rộng.

\begin{table}[H]
\centering
\caption{Phân chia nội dung thực hành theo hai khối yêu cầu}
\label{tab:split}
\small
\begin{tabularx}{\linewidth}{@{}l l X l@{}}
\toprule
\textbf{Khối} & \textbf{Mã câu} & \textbf{Nội dung} & \textbf{Tệp / CSDL} \\
\midrule
Bắt buộc & S01--S15 & Truy vấn cơ bản, CRUD, JOIN, tổng hợp, subquery, index & \code{sql/answers.sql} \\
         &         &                                                                 & \code{data/sales.db} \\
Bổ sung  & S16--S21 & Doanh thu theo tháng, khách chưa mua, sản phẩm chưa bán, & \code{sql/answers\_homework.sql} \\
(về nhà) &         & xếp hạng, phân tích kế hoạch, so sánh bảng tạm             & \code{data/sales\_extended.db} \\
Chứng minh index & S15, S20 & Cặp \code{EXPLAIN QUERY PLAN} trước/sau index & \code{sql/explain.sql} \\
\bottomrule
\end{tabularx}
\end{table}

% =====================================================================
\section{Môi trường thực hành và quy trình}
% =====================================================================

\subsection{Môi trường}

\begin{table}[H]
\centering
\caption{Môi trường thực hành}
\label{tab:env}
\small
\begin{tabularx}{\linewidth}{@{}l X@{}}
\toprule
\textbf{Thành phần} & \textbf{Giá trị} \\
\midrule
Nền tảng          & Google Cloud Shell (\code{https://shell.cloud.google.com}), chế độ có \code{\$HOME} bền vững \\
Ngôn ngữ          & Python 3.11.2 (thư viện chuẩn \code{sqlite3}, không cần cài thêm gói) \\
Hệ quản trị CSDL  & SQLite 3.40.1, chế độ tệp đơn \code{.db}, không cần máy chủ \\
Công cụ tương tác & \code{sql\_console.py} (console do đề cung cấp, có bổ sung tham số đường dẫn CSDL) \\
Công cụ dòng lệnh & \code{build\_db.py}, \code{generate\_data.py}, \code{run\_sql.py} \\
Chi phí           & 0 đồng --- không bật Billing, không tạo tài nguyên tính phí \\
\bottomrule
\end{tabularx}
\end{table}

Phiên bản SQLite ghi trong Bảng~\ref{tab:env} là phiên bản tại thời điểm kiểm chứng. Trên Cloud
Shell con số này có thể khác, nhưng toàn bộ câu lệnh của bài không phụ thuộc phiên bản; riêng
hai hàm cửa sổ \code{RANK()} và \code{DENSE\_RANK()} yêu cầu SQLite $\geq$ 3.25 (công bố 2018),
thấp hơn nhiều so với mọi bản Cloud Shell hiện hành.

\subsection{Quy trình thực hiện}

\begin{enumerate}[leftmargin=1.2cm,itemsep=2pt]
  \item Tạo cấu trúc bảng bằng \code{sql/schema.sql} (bật \code{PRAGMA foreign\_keys = ON});
  \item Nạp dữ liệu mẫu trong đề bằng \code{sql/seed.sql} và đối chiếu số dòng 5/5/6/8;
  \item Thực hiện 15 câu bắt buộc trên \code{data/sales.db};
  \item Sinh bộ dữ liệu mở rộng bằng \code{generate\_data.py} vào \code{data/sales\_extended.db};
  \item Thực hiện 6 câu bổ sung trên bộ dữ liệu mở rộng;
  \item Chạy \code{sql/explain.sql} để ghi lại kế hoạch thực thi trước và sau index.
\end{enumerate}

Một yêu cầu kỹ thuật quan trọng được đặt ra ngay từ đầu: \textbf{file đáp án phải chạy lại được
nhiều lần mà không làm sai dữ liệu}. Cụ thể, S04 thêm khách hàng thì S06 phải xoá đúng khách
hàng đó; S05 tăng giá 10\% chỉ được nhân một lần. Cách giải quyết trình bày ở mục~\ref{sec:s04s06}.

% =====================================================================
\section{Mô hình dữ liệu}
% =====================================================================

\subsection{Sơ đồ quan hệ}

Cơ sở dữ liệu gồm bốn bảng với ba quan hệ một--nhiều, đúng như đề bài quy định:

\begin{figure}[H]
\centering
\begin{tikzpicture}[
  entity/.style={draw, rounded corners=3pt, fill=blue!6, align=left,
                 font=\footnotesize\ttfamily, inner sep=5pt, text width=3.3cm},
  lnk/.style={-{Latex[length=2mm]}, thick, gray!70!black},
  lbl/.style={font=\scriptsize, fill=white, inner sep=1.5pt}
]
  \node[entity] (cust)  at (-4.6,0)  {customers\\ \textbf{customer\_id} PK\\ name\\ city};
  \node[entity] (ord)   at (0,0)     {orders\\ \textbf{order\_id} PK\\ customer\_id FK\\ order\_date};
  \node[entity] (item)  at (4.6,0)   {order\_items\\ \textbf{order\_id} PK,FK\\ \textbf{product\_id} PK,FK\\ quantity};
  \node[entity] (prod)  at (4.6,-2.6){products\\ \textbf{product\_id} PK\\ name, category\\ price};

  \draw[lnk] (cust.east) -- node[lbl,above]{1\,:\,N} (ord.west);
  \draw[lnk] (ord.east)  -- node[lbl,above]{1\,:\,N} (item.west);
  \draw[lnk] (prod.north) -- node[lbl,right]{1\,:\,N} (item.south);
\end{tikzpicture}
\caption{Sơ đồ quan hệ của cơ sở dữ liệu bán hàng}
\label{fig:erd}
\end{figure}

\subsection{Định nghĩa bảng}

\begin{lstlisting}[language=SQL,caption={Trích \code{sql/schema.sql}: bốn bảng và các ràng buộc}]
CREATE TABLE customers (
    customer_id INTEGER PRIMARY KEY,
    name        TEXT NOT NULL,
    city        TEXT
);

CREATE TABLE products (
    product_id INTEGER PRIMARY KEY,
    name       TEXT NOT NULL,
    category   TEXT,
    price      REAL CHECK(price >= 0)
);

CREATE TABLE orders (
    order_id    INTEGER PRIMARY KEY,
    customer_id INTEGER NOT NULL,
    order_date  TEXT,
    FOREIGN KEY(customer_id) REFERENCES customers(customer_id)
);

CREATE TABLE order_items (
    order_id   INTEGER,
    product_id INTEGER,
    quantity   INTEGER CHECK(quantity > 0),
    PRIMARY KEY(order_id, product_id),
    FOREIGN KEY(order_id) REFERENCES orders(order_id),
    FOREIGN KEY(product_id) REFERENCES products(product_id)
);
\end{lstlisting}

\subsection{Phân tích thiết kế}

\begin{itemize}[leftmargin=1.2cm,itemsep=2pt]
  \item \textbf{Khoá chính tổ hợp} \code{(order\_id, product\_id)} của \code{order\_items} ngăn
        việc cùng một sản phẩm xuất hiện hai lần trong một đơn; muốn mua thêm thì phải tăng
        \code{quantity}.
  \item \textbf{Ràng buộc \code{CHECK}} \code{price >= 0} và \code{quantity > 0} chặn dữ liệu vô
        nghĩa ngay tại tầng CSDL, không phụ thuộc vào việc chương trình ứng dụng có kiểm tra hay
        không.
  \item \textbf{Khoá ngoại} bảo đảm không thể tạo đơn cho khách hàng không tồn tại. SQLite
        \emph{không} bật kiểm tra khoá ngoại mặc định, nên mọi kết nối trong bài đều chạy
        \code{PRAGMA foreign\_keys = ON} ngay sau khi mở CSDL. Đây là điểm khác biệt so với
        SQL Server/MySQL và là lỗi rất dễ gặp.
  \item \textbf{\code{order\_date} kiểu TEXT} theo đúng đề bài. SQLite không có kiểu ngày riêng;
        chuỗi \code{'YYYY-MM-DD'} vẫn so sánh và sắp xếp đúng thứ tự thời gian, đồng thời dùng
        được với \code{strftime(\%Y-\%m, ...)} cho câu S16.
\end{itemize}

% =====================================================================
\section{Dữ liệu và kiểm chứng việc nạp dữ liệu}
% =====================================================================

\subsection{Dữ liệu mẫu theo đề}

Dữ liệu mẫu được nạp bằng \code{sql/seed.sql} với \code{INSERT OR IGNORE} để chạy lại không sinh
lỗi trùng khoá. Kết quả \code{python3 build\_db.py --force} in ra đối chiếu trực tiếp với kỳ vọng
của đề:

\begin{table}[H]
\centering
\caption{Đối chiếu số dòng sau khi nạp dữ liệu mẫu}
\label{tab:counts}
\small
\begin{tabular}{@{}l r r l@{}}
\toprule
\textbf{Bảng} & \textbf{Kỳ vọng} & \textbf{Thực tế} & \textbf{Kết luận} \\
\midrule
\code{customers}   & 5 & 5 & khớp \\
\code{products}    & 5 & 5 & khớp \\
\code{orders}      & 6 & 6 & khớp \\
\code{order\_items}& 8 & 8 & khớp \\
\bottomrule
\end{tabular}
\end{table}

\subsection{Kiểm tra tính toàn vẹn}

Sau khi nạp, câu \code{PRAGMA foreign\_key\_check} trả về danh sách rỗng, nghĩa là mọi
\code{customer\_id} trong \code{orders} và mọi \code{order\_id}/\code{product\_id} trong
\code{order\_items} đều trỏ tới bản ghi tồn tại. Trên bộ dữ liệu mở rộng (mục~\ref{sec:extended}),
phép kiểm tra này cũng cho kết quả \code{ok} với 1.500 dòng đơn hàng.

% =====================================================================
\section{Kết quả phần bắt buộc (S01--S15)}
% =====================================================================

Toàn bộ 15 câu nằm trong \code{sql/answers.sql}. Khi chạy bằng
\code{python3 run\_sql.py sql/answers.sql}, file thực thi 18 câu lệnh (gồm cả các câu kiểm tra
trung gian) với \textbf{0 lỗi}, và chạy lần thứ hai vẫn cho 0 lỗi cùng kết quả giống hệt.

\begin{table}[H]
\centering
\caption{Tổng hợp kỹ thuật của 15 câu bắt buộc}
\label{tab:s01s15}
\small
\begin{tabularx}{\linewidth}{@{}l X l@{}}
\toprule
\textbf{Mã} & \textbf{Kỹ thuật SQL} & \textbf{Số dòng} \\
\midrule
S01 & \code{SELECT} danh sách cột, \code{ORDER BY} & 5 \\
S02 & \code{WHERE} so sánh số & 2 \\
S03 & \code{ORDER BY ... DESC} & 5 \\
S04 & \code{INSERT INTO ... VALUES} & 1 (thêm mới) \\
S05 & \code{UPDATE ... SET price = price * 1.1} & 1 \\
S06 & \code{DELETE} có điều kiện \code{NOT EXISTS} & 1 \\
S07 & \code{INNER JOIN} hai bảng & 6 \\
S08 & \code{JOIN} ba bảng, \code{GROUP BY}, \code{SUM} & 6 \\
S09 & \code{LEFT JOIN} + \code{GROUP BY} & 5 \\
S10 & \code{HAVING SUM(...) > 5000000} & 3 \\
S11 & subquery vô hướng \code{> (SELECT AVG(price))} & 1 \\
S12 & \code{HAVING COUNT(order\_id) >= 2} & 1 \\
S13 & \code{ORDER BY} + \code{LIMIT 3} & 3 \\
S14 & \code{CREATE INDEX} trên \code{orders(customer\_id)} & --- \\
S15 & \code{EXPLAIN QUERY PLAN} & 1 (kế hoạch) \\
\bottomrule
\end{tabularx}
\end{table}

\subsection{Truy vấn đọc cơ bản (S01--S03)}

Kết quả S01--S03 in ra đúng dữ liệu mẫu của đề. Ở S02, chỉ Laptop (15.000.000) và Monitor
(3.500.000) vượt ngưỡng 1 triệu; S03 sắp xếp giảm dần theo giá nên Laptop đứng đầu và Mouse đứng
cuối.

\begin{table}[H]
\centering
\caption{Kết quả S03 --- sản phẩm sắp xếp theo giá giảm dần (trước khi S05 đổi giá Mouse)}
\label{tab:s03}
\small
\begin{tabular}{@{}r l l r@{}}
\toprule
\textbf{product\_id} & \textbf{name} & \textbf{category} & \textbf{price} \\
\midrule
101 & Laptop     & Electronics & 15.000.000 \\
104 & Monitor    & Electronics &  3.500.000 \\
105 & Headphones & Accessories &    800.000 \\
103 & Keyboard   & Accessories &    650.000 \\
102 & Mouse      & Accessories &    250.000 \\
\bottomrule
\end{tabular}
\end{table}

\subsection{CRUD và tính lặp an toàn (S04--S06)}
\label{sec:s04s06}

Ba câu này là thao tác \emph{có ghi}, nên nếu viết ngây thơ thì mỗi lần chạy lại file sẽ sai dữ
liệu. Hai biện pháp được áp dụng:

\begin{lstlisting}[language=SQL,caption={S04--S06: thao tác ghi an toàn khi chạy lặp}]
-- S04: them khach hang moi (chay lai khong loi khoa chinh)
INSERT INTO customers (customer_id, name, city)
VALUES (6, 'Lan', 'Hue')
ON CONFLICT(customer_id) DO NOTHING;

-- S05: tang gia Mouse 10%, chi khi con o gia goc 250000
UPDATE products
SET price = ROUND(price * 1.1, 0)
WHERE product_id = 102 AND price = 250000;

-- S06: xoa dung khach vua them va chi khi khach do chua co don
DELETE FROM customers
WHERE customer_id = 6
  AND NOT EXISTS (SELECT 1 FROM orders WHERE orders.customer_id = customers.customer_id);
\end{lstlisting}

Điểm mấu chốt là điều kiện \code{AND price = 250000}: lần chạy thứ hai, giá đã là 275.000 nên
câu \code{UPDATE} không khớp dòng nào và trả về \code{Rows affected: 0}. Kết quả kiểm chứng sau
hai lần chạy liên tiếp: giá Mouse vẫn là \textbf{275.000}, số khách hàng trở về \textbf{5}. Điều
kiện \code{NOT EXISTS} ở S06 cài đặt đúng yêu cầu "khách vừa thêm \emph{chưa có đơn}" của đề và
đồng thời bảo vệ khoá ngoại: nếu khách đã có đơn, câu lệnh sẽ không xoá gì thay vì gây lỗi.

\begin{table}[H]
\centering
\caption{Trạng thái dữ liệu trước và sau khi thực hiện S04--S06}
\label{tab:crud}
\small
\begin{tabular}{@{}l l l@{}}
\toprule
\textbf{Thời điểm} & \textbf{Số khách hàng} & \textbf{Giá Mouse} \\
\midrule
Ban đầu (sau seed)          & 5 & 250.000 \\
Sau S04 \code{INSERT}       & 6 & 250.000 \\
Sau S05 \code{UPDATE}       & 6 & 275.000 \\
Sau S06 \code{DELETE}       & 5 & 275.000 \\
Chạy lại toàn bộ file lần 2 & 5 & 275.000 \\
\bottomrule
\end{tabular}
\end{table}

\subsection{JOIN và tổng hợp (S07--S10)}

S08 là trung tâm của phần này: doanh thu một đơn bằng tổng của
\code{quantity}~$\times$~\code{price} trên các dòng đơn của đơn đó, nên phải nối ba bảng
\code{orders} -- \code{order\_items} -- \code{products}.

\begin{lstlisting}[language=SQL,caption={S08 --- tổng tiền từng đơn hàng}]
SELECT o.order_id,
       SUM(i.quantity * p.price) AS total
FROM orders o
JOIN order_items i ON i.order_id = o.order_id
JOIN products p    ON p.product_id = i.product_id
GROUP BY o.order_id
ORDER BY total DESC;
\end{lstlisting}

\begin{table}[H]
\centering
\caption{Kết quả S08 (tổng tiền từng đơn) và S09 (doanh thu từng khách)}
\label{tab:s08s09}
\small
\begin{minipage}[t]{0.47\linewidth}
\centering
\begin{tabular}{@{}r r@{}}
\toprule
\textbf{order\_id} & \textbf{total} \\
\midrule
1001 & 15.550.000 \\
1004 & 15.000.000 \\
1006 &  7.000.000 \\
1002 &  4.150.000 \\
1003 &  1.600.000 \\
1005 &    825.000 \\
\midrule
Tổng & 44.125.000 \\
\bottomrule
\end{tabular}
\end{minipage}\hfill
\begin{minipage}[t]{0.47\linewidth}
\centering
\begin{tabular}{@{}l r@{}}
\toprule
\textbf{Khách hàng} & \textbf{revenue} \\
\midrule
An (id 1)   & 17.150.000 \\
Chi (id 3)  & 15.000.000 \\
Ha (id 5)   &  7.000.000 \\
Binh (id 2) &  4.150.000 \\
Dung (id 4) &    825.000 \\
\bottomrule
\end{tabular}
\end{minipage}
\end{table}

Kiểm tra chéo bằng tay: đơn 1001 gồm 1 Laptop (15.000.000) và 2 Mouse (2 $\times$ 275.000), tức
15.550.000 --- khớp đúng kết quả máy trả về. Con số này cũng cho thấy \textbf{thứ tự thực hiện có
ảnh hưởng đến kết quả}: vì S05 đã tăng giá Mouse trước khi chạy S08, tổng tiền dùng giá
275.000 chứ không phải 250.000. Báo cáo ghi rõ điểm này để người đọc không nhầm là sai số.

S10 lọc theo kết quả tổng hợp nên phải dùng \code{HAVING} chứ không phải \code{WHERE}: chỉ ba
khách An, Chi, Ha vượt mốc 5 triệu.

\subsection{Subquery, HAVING theo số đơn, và TOP-N (S11--S13)}

\begin{itemize}[leftmargin=1.2cm,itemsep=2pt]
  \item \textbf{S11} so sánh với subquery vô hướng \code{(SELECT AVG(price) FROM products)}.
        Sau khi S05 đổi giá, giá trung bình là 4.045.000, nên chỉ Laptop (15.000.000) vượt trung
        bình. Subquery này \emph{độc lập} với truy vấn ngoài (không tương quan), SQLite chỉ cần
        tính một lần.
  \item \textbf{S12} đếm số đơn của từng khách rồi lọc bằng \code{HAVING COUNT(o.order\_id) >= 2};
        kết quả đúng một khách: An với 2 đơn (1001, 1003).
  \item \textbf{S13} dùng \code{ORDER BY revenue DESC LIMIT 3} để lấy top 3 sản phẩm theo doanh
        thu: Laptop 30.000.000 (bán trong 1001 và 1004), Monitor 10.500.000, Headphones
        1.600.000.
\end{itemize}

Lưu ý kỹ thuật: \code{LIMIT} không xử lý trường hợp "đồng hạng ở vị trí thứ 3" (nếu hai sản phẩm
cùng doanh thu bằng sản phẩm thứ ba, \code{LIMIT 3} chỉ trả về một trong hai). Muốn xử lý đúng
phải dùng \code{RANK()} như ở S19; đây là lý do S19 được xem là bước nâng cấp tự nhiên của S13.

% =====================================================================
\section{Index và phân tích kế hoạch thực thi (S14--S15)}
% =====================================================================

\subsection{Cách thực hiện}

Câu S15 yêu cầu so sánh kế hoạch \emph{trước} và \emph{sau} khi có index, nên không thể chỉ chạy
một lần sau khi đã tạo index. File \code{sql/explain.sql} được viết để tự đưa CSDL về trạng thái
cần thiết: \code{DROP INDEX IF EXISTS} $\rightarrow$ \code{EXPLAIN} $\rightarrow$
\code{CREATE INDEX IF NOT EXISTS} $\rightarrow$ \code{EXPLAIN}. Nhờ vậy file chạy đúng trên bất kỳ
bản sao nào của CSDL, ở bất kỳ trạng thái nào.

\begin{lstlisting}[language=SQL,caption={Cặp EXPLAIN trước/sau index trong \code{sql/explain.sql}}]
DROP INDEX IF EXISTS idx_orders_customer;
EXPLAIN QUERY PLAN
SELECT order_id FROM orders WHERE customer_id = 1;

CREATE INDEX IF NOT EXISTS idx_orders_customer ON orders(customer_id);

EXPLAIN QUERY PLAN
SELECT order_id FROM orders WHERE customer_id = 1;
\end{lstlisting}

\subsection{Kết quả đo được}

\begin{table}[H]
\centering
\caption{Kế hoạch thực thi trước và sau khi tạo index \code{idx\_orders\_customer}}
\label{tab:explain}
\small
\begin{tabularx}{\linewidth}{@{}l X@{}}
\toprule
\textbf{Thời điểm} & \textbf{\code{detail} do EXPLAIN QUERY PLAN trả về} \\
\midrule
Trước index & \code{SCAN orders} \\
Sau index   & \code{SEARCH orders USING COVERING INDEX idx\_orders\_customer (customer\_id=?)} \\
\bottomrule
\end{tabularx}
\end{table}

\subsection{Giải thích}

\begin{itemize}[leftmargin=1.2cm,itemsep=3pt]
  \item \code{SCAN orders} nghĩa là SQLite đọc tuần tự \emph{toàn bộ} bảng rồi mới lọc
        \code{customer\_id = 1}. Với 6 dòng thì không đáng kể, nhưng chi phí tăng tuyến tính theo
        số dòng: 500 đơn là 500 lần đọc, 1 triệu đơn là 1 triệu lần đọc.
  \item \code{SEARCH} nghĩa là đi xuống cây B-tree của index để tìm trực tiếp các dòng có
        \code{customer\_id = 1}. Chi phí tỉ lệ với $\log n$ cộng số dòng khớp, không phụ thuộc
        toàn bảng.
  \item \code{COVERING INDEX} là chi tiết đáng chú ý: câu truy vấn chỉ cần \code{order\_id},
        mà index trên \code{customer\_id} của bảng có \code{INTEGER PRIMARY KEY} đã chứa sẵn
        rowid (chính là \code{order\_id}). Vì vậy SQLite trả lời \emph{hoàn toàn từ index}, không
        phải quay lại bảng --- nhanh hơn cả \code{SEARCH ... USING INDEX}. Nếu truy vấn lấy thêm
        \code{order\_date} (như câu S20), kế hoạch sẽ là \code{SEARCH orders USING INDEX}, tức
        phải đọc thêm bảng.
  \item Khi dữ liệu lớn hơn, câu S20 cho thấy index còn giúp cả truy vấn tổng hợp:
        \code{SEARCH o USING COVERING INDEX idx\_orders\_customer (customer\_id>?)} là bước lọc
        \code{orders}; các bước còn lại đã dùng khoá chính/khoá tổ hợp
        (\code{sqlite\_autoindex\_order\_items\_1}, \code{INTEGER PRIMARY KEY}), nên không cần
        thêm index nào nữa.
\end{itemize}

\subsection{Vì sao không kết luận bằng thời gian chạy}

Với 6 đơn hàng, chênh lệch thời gian giữa \code{SCAN} và \code{SEARCH} nằm dưới ngưỡng đo của
đồng hồ hệ thống, nên nếu so sánh bằng \code{wall-clock} sẽ kết luận sai rằng index "không có
tác dụng". Cách đúng là đọc \emph{kế hoạch} (số bước, kiểu truy cập, có phải sort thêm hay
không), và khi cần số liệu định lượng thì dùng \code{EXPLAIN QUERY PLAN} kết hợp đo trên dữ liệu
đủ lớn. Đây cũng là điểm đề bài nhấn mạnh: dữ liệu rất nhỏ không đủ để suy luận index luôn giúp
nhanh hơn.

% =====================================================================
\section{Kết quả phần bổ sung (S16--S21)}
\label{sec:extended}
% =====================================================================

\subsection{Bộ dữ liệu mở rộng}

\code{generate\_data.py} sinh dữ liệu với \code{random.seed(2026)} để tái lập được, ghi vào
\code{data/sales\_extended.db} và \textbf{không} ghi đè \code{sales.db}. Kịch bản được thiết kế
có chủ đích để các câu S17, S18 có kết quả: khoảng 8\% khách hàng không phát sinh đơn, và bốn
sản phẩm cuối danh mục không xuất hiện trong bất kỳ đơn hàng nào.

\begin{table}[H]
\centering
\caption{Quy mô bộ dữ liệu mở rộng so với yêu cầu của đề}
\label{tab:extended}
\small
\begin{tabular}{@{}l r r l@{}}
\toprule
\textbf{Đối tượng} & \textbf{Yêu cầu} & \textbf{Thực tế} & \textbf{Ghi chú} \\
\midrule
Khách hàng    & $\geq$ 100 & 100  & gồm 5 khách của dữ liệu mẫu \\
Sản phẩm      & $\geq$ 30  & 30   & gồm 5 sản phẩm của dữ liệu mẫu \\
Đơn hàng      & $\geq$ 500 & 500  & id 1001--1500 \\
Dòng đơn hàng & $\geq$ 1.000 & 1.500 & mỗi đơn 1--5 sản phẩm \\
Khoá ngoại hợp lệ & bắt buộc & \code{ok} & \code{PRAGMA foreign\_key\_check} rỗng \\
\bottomrule
\end{tabular}
\end{table}

\subsection{S16 --- Doanh thu theo tháng}

Truy vấn nhóm theo \code{strftime('\%Y-\%m', o.order\_date)}. Dữ liệu trải 12 tháng (10/2025 --
10/2026) nên nhóm cuối chỉ có 7 đơn --- hiện tượng \emph{tháng dở dang} cần lưu ý khi đọc báo
cáo doanh thu thực tế.

\begin{table}[H]
\centering
\caption{S16 --- doanh thu và số đơn theo tháng trên bộ dữ liệu mở rộng}
\label{tab:s16}
\small
\begin{tabular}{@{}l r r@{}}
\toprule
\textbf{Tháng} & \textbf{Số đơn} & \textbf{Doanh thu (đ)} \\
\midrule
2025-10 & 53 & 1.142.485.000 \\
2025-11 & 39 &   978.725.000 \\
2025-12 & 36 &   797.655.000 \\
2026-01 & 38 &   780.995.000 \\
2026-02 & 37 &   870.780.000 \\
2026-03 & 46 & 1.051.310.000 \\
2026-04 & 47 &   947.890.000 \\
2026-05 & 42 &   838.505.000 \\
2026-06 & 35 &   813.970.000 \\
2026-07 & 40 &   696.115.000 \\
2026-08 & 44 &   899.005.000 \\
2026-09 & 36 &   853.310.000 \\
2026-10 &  7 &    48.400.000 \\
\bottomrule
\end{tabular}
\end{table}

\subsection{S17, S18 --- Khách chưa mua và sản phẩm chưa bán}

Hai câu này đều là bài toán "tìm phần tử không có bản ghi tương ứng", nhưng dùng hai kỹ thuật
khác nhau để so sánh: S17 dùng \code{LEFT JOIN ... WHERE ... IS NULL}, S18 dùng
\code{NOT EXISTS}. Cả hai đều đúng và có kế hoạch thực thi tương đương trên SQLite; \code{NOT EXISTS} thường được ưu tiên vì diễn đạt trực tiếp ngữ nghĩa và không sinh cột \code{NULL} phụ.

\begin{lstlisting}[language=SQL,caption={S17 và S18}]
-- S17: khach hang chua co don hang nao
SELECT c.customer_id, c.name, c.city
FROM customers c
LEFT JOIN orders o ON o.customer_id = c.customer_id
WHERE o.order_id IS NULL
ORDER BY c.customer_id;

-- S18: san pham chua tung ban duoc
SELECT p.product_id, p.name, p.category, p.price
FROM products p
WHERE NOT EXISTS (
    SELECT 1 FROM order_items i WHERE i.product_id = p.product_id
)
ORDER BY p.product_id;
\end{lstlisting}

\begin{table}[H]
\centering
\caption{Kết quả S17 (khách chưa có đơn) và S18 (sản phẩm chưa bán) trên dữ liệu mở rộng}
\label{tab:s17s18}
\small
\begin{minipage}[t]{0.44\linewidth}
\centering
\begin{tabular}{@{}r l l@{}}
\toprule
\textbf{id} & \textbf{name} & \textbf{city} \\
\midrule
37 & Ho Vy        & Cantho \\
38 & Bui Son      & Nhatrang \\
43 & Do Oanh      & HCMC \\
56 & Dang Lan     & Nhatrang \\
83 & Le An        & HCMC \\
89 & Le Dung      & Cantho \\
92 & Le Dung      & HCMC \\
94 & Hoang Quang  & Hanoi \\
\bottomrule
\end{tabular}
\\[2pt]{\footnotesize 8 khách hàng}
\end{minipage}\hfill
\begin{minipage}[t]{0.5\linewidth}
\centering
\begin{tabular}{@{}r l r@{}}
\toprule
\textbf{id} & \textbf{name} & \textbf{price} \\
\midrule
127 & Graphics Card & 7.200.000 \\
128 & RAM 16GB      & 1.450.000 \\
129 & CPU Cooler    &   722.000 \\
130 & Power Supply  & 1.980.000 \\
\bottomrule
\end{tabular}
\\[2pt]{\footnotesize 4 sản phẩm}
\end{minipage}
\end{table}

\subsection{S19 --- Xếp hạng khách hàng bằng window function}

\begin{lstlisting}[language=SQL,caption={S19 --- RANK và DENSE\_RANK}]
SELECT c.customer_id,
       c.name,
       ROUND(SUM(i.quantity * p.price), 0) AS revenue,
       RANK()       OVER (ORDER BY SUM(i.quantity * p.price) DESC) AS rank,
       DENSE_RANK() OVER (ORDER BY SUM(i.quantity * p.price) DESC) AS dense_rank
FROM customers c
JOIN orders o      ON o.customer_id = c.customer_id
JOIN order_items i ON i.order_id = o.order_id
JOIN products p    ON p.product_id = i.product_id
GROUP BY c.customer_id, c.name
ORDER BY revenue DESC;
\end{lstlisting}

\begin{table}[H]
\centering
\caption{S19 --- top 10 khách hàng theo doanh thu}
\label{tab:s19}
\small
\begin{tabular}{@{}r l r r r@{}}
\toprule
\textbf{id} & \textbf{Khách hàng} & \textbf{Doanh thu (đ)} & \textbf{rank} & \textbf{dense\_rank} \\
\midrule
 5 & Ha            & 314.565.000 &  1 &  1 \\
52 & Ho Mai        & 281.910.000 &  2 &  2 \\
95 & Hoang Chi     & 267.400.000 &  3 &  3 \\
39 & Hoang Yen     & 256.445.000 &  4 &  4 \\
36 & Hoang An      & 229.010.000 &  5 &  5 \\
85 & Ho Linh       & 220.360.000 &  6 &  6 \\
88 & Ho Oanh       & 202.410.000 &  7 &  7 \\
81 & Pham An       & 199.155.000 &  8 &  8 \\
79 & Tran Khanh    & 187.115.000 &  9 &  9 \\
 1 & An            & 186.680.000 & 10 & 10 \\
\bottomrule
\end{tabular}
\end{table}

Với bộ dữ liệu này không có hai khách nào trùng doanh thu nên \code{rank} và \code{dense\_rank}
trùng nhau. Khác biệt chỉ lộ ra khi có đồng hạng: nếu hai khách cùng đứng thứ 2, \code{RANK()}
gán 2, 2 rồi \emph{nhảy} sang 4 (bỏ 3), còn \code{DENSE\_RANK()} gán 2, 2 rồi 3. Vì vậy khi báo
cáo "top 10 khách hàng" nên dùng \code{DENSE\_RANK} để không bị thiếu người, còn khi cần biết
chính xác có bao nhiêu người đứng trên một khách thì dùng \code{RANK}.

\subsection{S20 --- Kế hoạch thực thi sau khi mở rộng dữ liệu}

\begin{table}[H]
\centering
\caption{S20 --- so sánh kế hoạch thực thi giữa dữ liệu mẫu và dữ liệu mở rộng}
\label{tab:s20}
\small
\begin{tabularx}{\linewidth}{@{}l l X@{}}
\toprule
\textbf{Dữ liệu} & \textbf{Index} & \textbf{Kế hoạch} \\
\midrule
6 đơn (mẫu)   & chưa có & \code{SCAN orders} \\
6 đơn (mẫu)   & có      & \code{SEARCH orders USING COVERING INDEX idx\_orders\_customer (customer\_id=?)} \\
500 đơn       & chưa có & \code{SCAN orders} \\
500 đơn       & có      & \code{SEARCH orders USING INDEX idx\_orders\_customer (customer\_id=?)} \\
              &         & \code{USE TEMP B-TREE FOR ORDER BY} (do có \code{ORDER BY order\_date}) \\
\bottomrule
\end{tabularx}
\end{table}

Hai điểm rút ra. Thứ nhất, \emph{kiểu truy cập} đổi từ đọc toàn bảng sang tìm theo index ngay cả
khi dữ liệu chỉ 500 dòng --- đây là bằng chứng định tính, không phụ thuộc tốc độ máy. Thứ hai,
khi truy vấn cần thêm cột không nằm trong index (\code{order\_date}) và có \code{ORDER BY} trên
cột đó, SQLite phải dựng B-tree tạm để sắp xếp; nếu truy vấn này chạy thường xuyên, index tổ hợp
\code{(customer\_id, order\_date)} sẽ loại bỏ được bước sort.

\subsection{S21 --- SQL Server (\#table / \#\#table / tempdb) và SQLite TEMP TABLE}

\begin{table}[H]
\centering
\caption{So sánh bảng tạm giữa SQL Server và SQLite}
\label{tab:s21}
\small
\begin{tabularx}{\linewidth}{@{}l X l@{}}
\toprule
\textbf{Tiêu chí} & \textbf{SQL Server} & \textbf{SQLite} \\
\midrule
Bảng tạm cục bộ &
\code{\#t}: chỉ session (hoặc procedure) tạo ra nó nhìn thấy; tự xoá khi session/procedure kết
thúc; nằm trong \code{tempdb} &
\code{CREATE TEMP TABLE t}: chỉ connection tạo ra nó nhìn thấy; tự mất khi đóng kết nối \\
Bảng tạm toàn cục &
\code{\#\#t}: mọi session đều thấy; mất khi session cuối cùng đang tham chiếu đóng lại &
Không có khái niệm tương đương \\
Nơi lưu trữ &
\code{tempdb} --- CSDL hệ thống dùng chung, cũng là nơi chứa biến bảng, sort/hash spill;
truy vấn trực tiếp được \code{tempdb..\#t} &
Schema \code{temp} gắn với kết nối; có thể ở RAM hoặc tệp tạm tuỳ \code{PRAGMA temp\_store}
(0 mặc định, 1 tệp, 2 bộ nhớ) \\
Phạm vi tên &
Tên \code{\#}/\code{\#\#} phân biệt bằng tiền tố; bảng thường trùng tên vẫn phân biệt được &
Bảng trong \code{temp} được tìm \emph{trước} \code{main}, nên trùng tên sẽ \emph{che} bảng
thường cho tới khi đóng kết nối \\
Cú pháp xoá &
\code{DROP TABLE \#t} &
\code{DROP TABLE temp.t} \\
\bottomrule
\end{tabularx}
\end{table}

Điểm dễ nhầm: SQLite \emph{không} có bảng tạm toàn cục. Khi cần chia sẻ dữ liệu trung gian giữa
nhiều tiến trình, phải dùng bảng thường trong tệp CSDL (hoặc bảng \code{ATTACH}), chứ không thể
tái tạo hành vi của \code{\#\#table}.

Phần minh hoạ bằng \code{TEMP TABLE} đã chạy trong \code{sql/answers\_homework.sql}: tạo bảng
tạm tổng hợp doanh thu theo tháng, truy vấn 3 tháng cao nhất (kết quả: 2025-10, 2026-03,
2025-11), rồi \code{DROP TABLE temp.revenue\_by\_month}.

% =====================================================================
\section{Thảo luận và bài học}
% =====================================================================

\begin{enumerate}[leftmargin=1.2cm,itemsep=4pt]
  \item \textbf{Khoá ngoại của SQLite là tuỳ chọn.} Mặc định SQLite không kiểm tra khoá ngoại;
        thiếu \code{PRAGMA foreign\_keys = ON} thì ràng buộc trong \code{CREATE TABLE} chỉ là
        chú thích. Bài làm bật pragma này ở mọi kết nối (\code{sql\_console.py},
        \code{build\_db.py}, \code{run\_sql.py}, \code{generate\_data.py}).
  \item \textbf{Script có ghi phải lặp an toàn.} Dùng \code{INSERT ... ON CONFLICT DO NOTHING},
        điều kiện chặt trong \code{UPDATE}, và điều kiện \code{NOT EXISTS} trong \code{DELETE} là
        ba mẫu nhỏ nhưng quyết định chất lượng file nộp: chạy lại không sai dữ liệu, người chấm
        không phải dựng lại CSDL thủ công.
  \item \textbf{Index không miễn phí.} Mỗi index làm \code{INSERT}/\code{UPDATE}/\code{DELETE}
        chậm thêm và tốn dung lượng, nên chỉ tạo index cho cột thường xuyên lọc/nối
        (\code{orders.customer\_id}) chứ không tạo tràn lan.
  \item \textbf{Đọc kế hoạch thay vì đo thời gian.} \code{EXPLAIN QUERY PLAN} cho biết
        \emph{cách} CSDL truy cập dữ liệu; đó mới là thông tin ổn định giữa các máy. Đặc biệt chú
        ý \code{COVERING INDEX} và \code{USE TEMP B-TREE FOR ORDER BY} vì chúng chỉ ra chi phí ẩn.
  \item \textbf{Dữ liệu quyết định kết luận.} Cùng một câu S08 cho tổng tiền khác nhau tuỳ thời
        điểm chạy S05; cùng một truy vấn cho kế hoạch khác nhau giữa 6 và 500 dòng. Báo cáo phải
        ghi rõ trạng thái dữ liệu khi lấy số liệu, nếu không sẽ mâu thuẫn khi người chấm chạy lại.
  \item \textbf{Thứ tự thực thi có ngữ nghĩa.} S16 đã cho thấy tháng cuối bị dở dang; nếu báo
        cáo doanh thu theo tháng thì cần nêu rõ kỳ dữ liệu kết thúc ngày nào, tránh kết luận
        "doanh thu tháng 10 giảm mạnh".
\end{enumerate}

% =====================================================================
\section{Kết luận}
% =====================================================================

Bài LAB 01 đã hoàn thành đầy đủ 21 yêu cầu. Cơ sở dữ liệu bốn bảng được tạo đúng thiết kế của đề
với khoá chính, khoá ngoại và ràng buộc \code{CHECK}; dữ liệu mẫu nạp đúng 5/5/6/8 dòng; toàn bộ
câu bắt buộc S01--S15 và câu bổ sung S16--S21 chạy không lỗi, chạy lặp vẫn cho kết quả ổn định.
Phần index chứng minh được bằng kế hoạch thực thi rằng việc tạo index chuyển truy vấn từ
\code{SCAN} sang \code{SEARCH} (thậm chí \code{COVERING INDEX}), đồng thời chỉ ra giới hạn của
việc suy luận hiệu năng trên dữ liệu quá nhỏ.

Hướng mở rộng tiếp theo: thêm view/báo cáo doanh thu, thử nghiệm index tổ hợp
\code{(customer\_id, order\_date)} cho truy vấn S20, và so sánh với phiên bản MongoDB ở LAB 02
để làm rõ trade-off giữa mô hình chuẩn hoá và mô hình nhúng.

% =====================================================================
\ifwithappendix
\clearpage
\appendix
% =====================================================================

\section{Cấu trúc bài nộp và lệnh tái lập}

\begin{lstlisting}[caption={Cấu trúc thư mục \code{lab01\_sql/}}]
lab01_sql/
  sql_console.py          console SQLite tuong tac (bat buoc nop)
  build_db.py             dung data/sales.db tu sql/schema.sql + sql/seed.sql
  generate_data.py        sinh data/sales_extended.db (100/30/500/1500)
  run_sql.py              chay file .sql bang Python sqlite3
  sql/
    schema.sql            4 bang + khoa ngoai + CHECK
    seed.sql              du lieu mau dung de bai (5/5/6/8)
    answers.sql           BAT BUOC  - S01..S15
    answers_homework.sql  BO SUNG  - S16..S21
    explain.sql           EXPLAIN QUERY PLAN truoc/sau index
  data/                   sales.db, sales_extended.db (sinh tu dong)
  report/                 lab01_report.tex + PDF
  evidence/               anh minh chung
\end{lstlisting}

\begin{lstlisting}[language=bash,caption={Tái lập toàn bộ kết quả trong báo cáo}]
cd lab01_sql

# 1) Du lieu mau trong de -> ky vong 5 / 5 / 6 / 8
python3 build_db.py --force

# 2) Phan BAT BUOC (S01-S15) tren data/sales.db
python3 run_sql.py sql/answers.sql

# 3) Du lieu mo rong + phan BO SUNG (S16-S21)
python3 generate_data.py --force
python3 run_sql.py sql/answers_homework.sql --db data/sales_extended.db

# 4) Chung minh index: SCAN orders -> SEARCH ... USING INDEX
python3 run_sql.py sql/explain.sql
python3 run_sql.py sql/explain.sql --db data/sales_extended.db
\end{lstlisting}

\section{Danh mục ảnh minh chứng cần chụp}

Các ảnh dưới đây dán vào Phụ lục này (hoặc nộp kèm trong thư mục \code{evidence/}) sau khi thực
hiện trên Google Cloud Shell. Tên tệp theo đúng quy ước để khớp với thư mục \code{evidence/}.

\begin{table}[H]
\centering
\caption{Danh mục ảnh minh chứng}
\label{tab:evidence}
\small
\begin{tabularx}{\linewidth}{@{}l X@{}}
\toprule
\textbf{Tệp} & \textbf{Nội dung cần thấy} \\
\midrule
\code{01\_env.png} & \code{whoami}, \code{pwd}, \code{python3 --version}, \code{sqlite3 --version} trên Cloud Shell \\
\code{02\_tables\_count.png} & \code{SELECT COUNT(*)} của 4 bảng: 5 / 5 / 6 / 8 \\
\code{03\_console.png} & \code{sql\_console.py} đang chạy, một câu lệnh nhiều dòng và \code{EXIT} \\
\code{04\_s08\_result.png} & Kết quả S08 (bảng tổng tiền 6 đơn) \\
\code{05\_explain\_before\_after.png} & Hai lần \code{EXPLAIN QUERY PLAN}: \code{SCAN} rồi \code{SEARCH ... INDEX} \\
\code{06\_expanded\_counts.png} & Kết quả \code{generate\_data.py}: 100 / 30 / 500 / 1500 \\
\code{07\_s19\_rank.png} & Kết quả S19 với cột \code{rank} và \code{dense\_rank} \\
\bottomrule
\end{tabularx}
\end{table}

\section{Mã nguồn phần bắt buộc (trích lõi truy vấn)}

Toàn bộ nội dung tệp \code{sql/answers.sql} nằm trong repo; phần dưới trích các truy vấn chính
để đối chiếu nhanh với báo cáo. Chú thích tiếng Việt trong tệp gốc được lược bỏ ở đây để bản in
gọn hơn.

\begin{lstlisting}[language=SQL,caption={S01--S13: truy vấn chính}]
-- S01
SELECT name, city FROM customers ORDER BY customer_id;

-- S02
SELECT product_id, name, price FROM products
WHERE price > 1000000 ORDER BY price DESC;

-- S03
SELECT product_id, name, category, price FROM products ORDER BY price DESC;

-- S07
SELECT o.order_id, c.name AS customer_name, o.order_date
FROM orders o JOIN customers c ON c.customer_id = o.customer_id
ORDER BY o.order_id;

-- S09
SELECT c.customer_id, c.name, SUM(i.quantity * p.price) AS revenue
FROM customers c
LEFT JOIN orders o      ON o.customer_id = c.customer_id
LEFT JOIN order_items i ON i.order_id = o.order_id
LEFT JOIN products p    ON p.product_id = i.product_id
GROUP BY c.customer_id, c.name ORDER BY revenue DESC;

-- S10
SELECT c.customer_id, c.name, SUM(i.quantity * p.price) AS revenue
FROM customers c
JOIN orders o      ON o.customer_id = c.customer_id
JOIN order_items i ON i.order_id = o.order_id
JOIN products p    ON p.product_id = i.product_id
GROUP BY c.customer_id, c.name
HAVING SUM(i.quantity * p.price) > 5000000
ORDER BY revenue DESC;

-- S11
SELECT product_id, name, price FROM products
WHERE price > (SELECT AVG(price) FROM products) ORDER BY price DESC;

-- S12
SELECT c.customer_id, c.name, COUNT(o.order_id) AS order_count
FROM customers c JOIN orders o ON o.customer_id = c.customer_id
GROUP BY c.customer_id, c.name HAVING COUNT(o.order_id) >= 2;

-- S13
SELECT p.product_id, p.name, SUM(i.quantity * p.price) AS revenue
FROM products p JOIN order_items i ON i.product_id = p.product_id
GROUP BY p.product_id, p.name ORDER BY revenue DESC LIMIT 3;

-- S14
CREATE INDEX IF NOT EXISTS idx_orders_customer ON orders(customer_id);
\end{lstlisting}

\section{Bảng đối chiếu yêu cầu của đề}

\begin{table}[H]
\centering
\caption{Đối chiếu tiêu chí chấm điểm với nội dung báo cáo}
\label{tab:rubric}
\small
\begin{tabularx}{\linewidth}{@{}X c l@{}}
\toprule
\textbf{Tiêu chí} & \textbf{Điểm} & \textbf{Vị trí trong báo cáo} \\
\midrule
Cài đặt và tạo dữ liệu & 2.0 & Mục 2, 3, 4; Bảng~\ref{tab:counts}, \ref{tab:env} \\
CRUD và câu cơ bản & 2.0 & Mục 5.1--5.3; Bảng~\ref{tab:s03}, \ref{tab:crud} \\
Truy vấn nâng cao & 3.0 & Mục 5.4--5.5; Bảng~\ref{tab:s08s09}; mục 8.2--8.5 \\
Index và giải thích & 1.5 & Mục 6; Bảng~\ref{tab:explain}, \ref{tab:s20} \\
Bài mở rộng, báo cáo & 1.5 & Mục 8; Bảng~\ref{tab:extended}, \ref{tab:s16}, \ref{tab:s19}, \ref{tab:s21} \\
\bottomrule
\end{tabularx}
\end{table}

\fi
% =====================================================================
%  TÀI LIỆU THAM KHẢO (luôn hiển thị)
% =====================================================================
\clearpage
\begin{thebibliography}{9}
\small
\bibitem{sqliteeqp} SQLite Consortium. \emph{EXPLAIN QUERY PLAN}.
  \url{https://www.sqlite.org/eqp.html}.
\bibitem{sqlitefk} SQLite Consortium. \emph{Foreign Key Support}.
  \url{https://www.sqlite.org/foreignkeys.html}.
\bibitem{sqlitetemp} SQLite Consortium. \emph{Temporary Files Used By SQLite} và
  \emph{CREATE TABLE (TEMP)}. \url{https://www.sqlite.org/tempfiles.html}.
\bibitem{pysqlite} Python Software Foundation. \emph{sqlite3 --- DB-API 2.0 interface for SQLite
  databases}. \url{https://docs.python.org/3/library/sqlite3.html}.
\bibitem{cloudshell} Google Cloud. \emph{Cloud Shell documentation}.
  \url{https://cloud.google.com/shell/docs}.
\bibitem{mssqltemp} Microsoft. \emph{Tables (Transact-SQL) --- temporary tables and tempdb}.
  \url{https://learn.microsoft.com/sql/t-sql/statements/create-table-transact-sql}.
\bibitem{sqlitewindow} SQLite Consortium. \emph{Window Functions}.
  \url{https://www.sqlite.org/windowfunctions.html}.
\end{thebibliography}

\end{document}
